\begin{studiumworked}{M22. Calentamiento electrotérmico de un gas académico · Intermedio}
\textbf{Enunciado.} Un calentador transfiere al gas el 65\,\% de 900 W eléctricos. Entran 1.2 g/s de gas académico a 300 K, con $c_p=2100$ J/(kg K) constante. Determine temperatura total de salida y cierre energía.

\textbf{Datos.} Potencia: $900\,\mathrm{W}$; Fracción absorbida: $0.65$; Caudal: $0.0012\,\mathrm{kg/s}$; Capacidad: $2100\,\mathrm{J\,kg^{-1}\,K^{-1}}$; Entrada: $300\,\mathrm{K}$.

\textbf{Incógnitas.} Magnitudes solicitadas y cierre físico.

\textbf{Hipótesis.} Modelo académico con hipótesis del capítulo. Datos originales; no representan hardware ni ensayo.

\textbf{Modelo físico.} Conservación y relaciones constitutivas declaradas.
\[\dot m c_p(T_{out}-T_{in})=\eta_hP_{el}\]
\[T_{out}=T_{in}+\eta_hP_{el}/(\dot m c_p)\]
\textbf{Desarrollo.}
\begin{enumerate}
\item Convertir 1.2 g/s a SI da $0.0012\,\mathrm{kg/s}$. La potencia al fluido es $0.65(900)=585$ W.
\item Despejar temperatura da $300+585/(0.0012\cdot2100)=532.1429$ K.
\item Los $315\,\mathrm{W}$ restantes pertenecen a otras rutas de energía de la frontera elegida. No se restan de nuevo los 585 W absorbidos.
\end{enumerate}
\textbf{Resultados.} $T_{out}=532.1429$ K.

\textbf{Interpretación.} El resultado corresponde a las hipótesis indicadas; su interpretación requiere conservar fronteras y unidades.

\textbf{Comprobación.} El aumento de entalpía por caudal recupera 585 W; la potencia absorbida es menor que la entrada.

\textbf{Límites.} No se modela descarga eléctrica, disociación ni tobera; no se deduce aún Isp.
\end{studiumworked}
